% Lösungen Einheit 5 von 5 – Kurvenuntersuchung im Sachzusammenhang (zusammenbau v0.1)
\einheitenkopf{Einheit 5 von 5 · Kurvenuntersuchung im Sachzusammenhang}

\erg{34}{a) stärkste Änderung: Extremum von $f'$, Wendepunkt – gefragt ist die größte Wachstumsrate. \quad b) Die Temperatur fällt von $90$ °C zunächst schnell, dann immer langsamer und nähert sich $20$ °C, der Raumtemperatur; sie sinkt nicht darunter. \quad c) Bei $t = 10$ steigt der Ballon am schnellsten; danach steigt er weiter, aber langsamer – der Wendepunkt ist kein Richtungswechsel. \quad d) $25$ cm – die steilste Stelle ist der Wendepunkt bei $t = 10$. \quad e) $f'(x) = -0{,}4x + 1{,}6 = 0$ für $x = 4$; $f(4) = 3{,}6$; Hochpunkt $H(4 | 3{,}6)$, der Bogen ist $3{,}6$ m hoch – nicht $4$ m, das ist die Stelle. \quad f) $W'(t) = -1{,}5t^2 + 12t$, $W'(8) = 0$; $W''(t) = -3t + 12$, $W''(8) = -12 < 0$; $W(8) = 228$ m³; Randwerte $W(0) = 100$, $W(10) = 200$ sind kleiner. \quad g) $c''(t) = 0 \Leftrightarrow 0{,}24t - 2{,}4 = 0 \Leftrightarrow t = 10$; $c(10) = 60\,\mathrm{e}^{-2} \approx 8{,}12$ mg pro Liter; $c'(t) = (6 - 1{,}2t)\,\mathrm{e}^{-0{,}2t}$, $c'(10) = -6\,\mathrm{e}^{-2} \approx -0{,}81$ mg pro Liter und Stunde. \quad h) $B'(t) = -3t^2 + 12t + 15 = -3(t - 5)(t + 1) = 0$ für $t = 5$ ($t = -1$ verwerfen); $B(5) = 200$; Randwerte $B(0) = 100$, $B(8) = 92$; die meisten Gäste, $200$, sind nach $5$ Stunden da. \quad i) $g(-3) = -3 + 4{,}5 = 1{,}5$: Höhe $2 \cdot 1{,}5 = 3$ m; die Halbkreise haben den Radius $1{,}5$ m, Breite $2 \cdot (3 + 1{,}5) = 9$ m. Wegen $3 \le 3{,}2$ und $9 \le 9{,}2$ passt die Bahn. Ohne Halbkreise wäre die Breite nur $6$ m. \quad j) $d(t)$ ist der Höhenunterschied, um den $A$ größer ist als $B$. Nach etwa $4{,}2$ Wochen ist dieser Unterschied am größten; er beträgt dann $d(4{,}2) \approx 2{,}6$ cm – eine Länge, keine Zeit.}
\erg{35}{a) $h(5) = 32$ bedeutet: $5$ Minuten nach dem Start fliegt der Drachen $32$ m hoch.}
\erg{36}{a) $18$ °C – die Mitte zwischen dem höchsten Wert $22$ °C und dem tiefsten Wert $14$ °C.}
\erg{37}{a) Fehler: Der Wendepunkt ist kein Richtungswechsel. Bei $t = 10$ steigt der Ballon am schnellsten, danach langsamer; sinken kann er erst nach dem Hochpunkt bei $t = 20$. \quad b) Der Zeitraum schneidet den Graphen ab. Steigt der Wasserstand bis zum Ende der Beobachtung, ist der letzte Wert der größte, auch ohne waagerechte Tangente – deshalb vergleicht man Hochpunkte immer mit den Randwerten. \quad c) $f'(x) = -0{,}012x^2 + 0{,}18x$, $f''(x) = -0{,}024x + 0{,}18 = 0$ für $x = 7{,}5$; $f(7{,}5) \approx 5{,}38$; $f'(7{,}5) = 0{,}675$; $\mathrm{tan}\,\alpha = 0{,}675$, also $\alpha \approx 34{,}0° > 30°$: die Bedingung ist nicht erfüllt. \quad d) Gerade durch $(0 | 0)$ und $(9 | 216)$; sie schneidet den Graphen von $K$ bei $x = 2$ und $x = 8$ und liegt dazwischen oberhalb: Gewinn für Mengen zwischen $2$ m³ und $8$ m³ – links von $2$ nicht, dort sind die Kosten höher.}

37d)

\begin{ksys}[xmin=0,xmax=9,ymin=0,ymax=300,xstep=1,ystep=50,klein] \funktionab{\x^3-9*\x^2+30*\x+16}{K}{0}{9} \gerade{24}{0}{E} \end{ksys}

